Complementar el análisis exploratorio mediante técnicas estadísticas y modelos predictivos que permitan identificar los factores con mayor capacidad explicativa sobre la mora y cuantificar su impacto en el riesgo crediticio.
Incorporar variables adicionales relacionadas con el perfil financiero del cliente, como historial de pagos, nivel de endeudamiento, relación cuota-ingreso, comportamiento crediticio previo y calificación de riesgo, con el fin de fortalecer los modelos de evaluación.
Dar seguimiento periódico a la evolución de la tasa de mora mediante indicadores y tableros de control que permitan detectar oportunamente cambios en el comportamiento de la cartera y facilitar la toma de decisiones.
Fortalecer las políticas de otorgamiento y seguimiento de créditos de acuerdo con las normas internas de gestión del riesgo crediticio, priorizando decisiones respaldadas por evidencia estadística y análisis de datos.